% Standalone appendix document - compiles on its own:
%   latexmk -pdf -outdir=build 03-agenda-2026-09-30.tex
\documentclass[a4paper,11pt]{article}
\usepackage[danish]{babel}
\usepackage{../../appendix-style}

\title{12.3 Dagsorden -- 30-09}
\author{SW4PRJ4 -- Group 3}
\date{30-09-2026}

\begin{document}
\maketitle

\section*{Mødeinformation}

\textbf{Dato:} 30-09-2026 (onsdag)

\textbf{Tid:}

\textbf{Sted:}

\section*{Dagsordenspunkter}

\begin{enumerate}
    \item \textbf{Sprint review -- sprint 1} (ca. 30 min)
    \begin{itemize}
        \item Alt, der står i \emph{Done}, er færdigt fra sprint 1: 12 kort (34 point). PMU-1, PMU-5, PMU-9, PMU-10, PMU-14, PMU-15, PMU-16, PMU-17, PMU-20, PMU-21, PMU-22 og PMU-23.
        \item Ikke færdigt: 2 kort (8 point) i \emph{In Progress} (PMU-2, PMU-12) og 6 kort (15 point) i \emph{To Do} (PMU-3, PMU-4, PMU-13, PMU-18, PMU-19, PMU-70). \emph{In Review} er tom.
        \item Den første spike er ikke helt i mål. Vi demonstrerer og tester den først, når der er noget, vi rent faktisk kan demonstrere fra ende til anden.
        \item Velocity for sprint 1 (fra Kaneo-boardet i dag): committed 57 point, done 34 point, spillover 23 point, scope 57 point. Tallene står i \texttt{docs/assets/data/sprint-velocity.csv}.
        \item Scope tæller kun estimerede kort. De 45 kort i backloggen (PMU-24 -- PMU-68) har endnu ingen point-label.
    \end{itemize}
    \item \textbf{Retrospective} (ca. 30 min)
    \begin{itemize}
        \item \textbf{Hvad gik godt?} AI-opsætningen og strukturen i hele projektet. Alle er sat ind i, hvordan opsætningen fungerer. Den første afgrænsning af projektet er så lille, at ingen bliver overvældet: man kan overskue helheden og dykke ned i sine egne opgaver.
        \item \textbf{Hvad gik skidt?} Man har ikke overblik over, hvad den anden undergruppe har lavet, før vi holder en demo, og den kan vi ikke holde endnu.
        \item \textbf{Hvad ændrer vi i sprint 2?}
        \begin{itemize}
            \item Hver sprint slutter med en demo, så alle også kender de dele, de ikke selv har arbejdet på.
            \item Et lille sæt overbliksdiagrammer (fx et E/R-diagram), som man kan slå op i for at forstå projektet. Kun diagrammer, der bliver brugt -- ingen ligegyldige diagrammer.
            \item Hver handling bliver et kort i Kaneo.
        \end{itemize}
        \item Skal definition of done justeres?
    \end{itemize}
    \item \textbf{Sprint planning -- sprint 2} (ca. 60 min)
    \begin{itemize}
        \item Kapacitet: hvor mange timer har hver af os i perioden 30/9 -- 13/10?
        \item Spillover: de 8 kort (23 point), som ikke er done, beholder label \emph{Sprint 1} og får \emph{Sprint 2} oveni, hvis de fortsætter. PMU-2 og PMU-12 er i gang. Skal PMU-3, PMU-4, PMU-13, PMU-18, PMU-19 og PMU-70 med i sprint 2 eller tilbage i backloggen?
        \item PMU-70 (sprint 1) og PMU-27 (sprint 2) handler begge om login med Keycloak i kunde-appen. Skal de slås sammen?
        \item Planning poker på sprint 2-kortene PMU-24 -- PMU-33. Ingen af dem har point-label eller ejer endnu.
        \item Kort til retrospective-handlingerne: demo ved sprintens afslutning og overbliksdiagrammer. E/R-diagrammet kan høre sammen med PMU-25 (EF-konfiguration og migration) og domænemodellen i PMU-2.
        \item Ready-kort fra backlog til \emph{To Do}: \texttt{Krav: FR-xx} som første linje, definition of done, point-label, epic-label og ejer.
        \item Sprint goal (én sætning, der nævner de krav sprinten leverer).
        \item Ejer og undergruppe på hvert kort, og opgaver for de første dage.
        \item Point i \emph{To Do} (committed).
    \end{itemize}
    \item \textbf{Eventuelt}
\end{enumerate}

\end{document}
